\begin{theorem}[7.2.5]
\label{IV.2.s7.2.5-fr}
Sea $A$ un anillo local noetheriano. Las condiciones siguientes son
equivalentes:
\begin{enumerate}
\item[\textup{a)}] Para todo anillo cociente íntegro $B$ de $A$, su completado
$\widehat B$ es estrictamente equidimensional \textup{(7.2.1)}.
\item[\textup{b)}] $A$ es formalmente catenario \textup{(7.1.9)} y las fibras del
morfismo canónico
$\operatorname{Spec}(\widehat A)\to\operatorname{Spec}(A)$
satisfacen $(\mathrm S_1)$ \textup{(es decir,} no tienen ningún ciclo primo
asociado inmerso\textup{)}.
\item[\textup{c)}] $A$ es universalmente catenario \textup{(5.6.2)} y, para todo
anillo cociente íntegro $B$ de $A$, el anillo $B^{(1)}$ es una
$B$-álgebra finita \textup{(cf. (7.2.4))}.
\item[\textup{d)}] $A$ es universalmente catenario y, para todo anillo
cociente íntegro $B$ de $A$ y todo anillo local $C=B_{\mathfrak q}$
en un ideal primo $\mathfrak q$ de $B$ tal que $\dim(C)\geqslant2$,
el anillo $C^{(\omega)}$ es una $C$-álgebra finita.
\item[\textup{e)}] $A$ es universalmente catenario y, para todo anillo
cociente íntegro $B$ de $A$ y todo anillo local $C=B_{\mathfrak q}$
en un ideal primo $\mathfrak q$ de $B$ tal que $\dim(C)\geqslant2$,
el anillo completado $\widehat C$ es tal que
$\operatorname{Spec}(\widehat C)$ no tiene ningún ciclo primo asociado
de dimensión 1.
\end{enumerate}
Además, cuando se cumplen estas condiciones, para todo anillo
cociente $B$ de $A$ que sea estrictamente equidimensional, su
completado $\widehat B$ es estrictamente equidimensional.
\end{theorem}
